\chapter{Revisão da literatura}
\label{cap:literatura}

\section{Linha de base da pesquisa}
\label{sec:baseline}

\subsection{Success Criteria e verificação \& validação}

Success Criteria verificáveis exigem desempenho mensurável, tolerâncias, escopo
temporal/espacial e coerência com o método de verificação escolhido~\cite{msfc3173,sis08methods}.
O gate ACCEPT/REJECT do ReqValLive operacionaliza a aprovação do critério
\emph{antes} da actividade de V\&V --- alinhado às oito considerações MSFC
(Performance, Environment, Restrictions, Checkpoints, entre outras).

\subsection{MBSE e SysML~v2}

O SysML~v2 oferece notação textual e APIs para modelos de requisitos~\cite{sysmlv2}.
Neste trabalho, o export SysML serve como \emph{esqueleto} para reuso MBSE; o laudo
HTML é a evidência primária de medição. O SysON actua como repositório do ciclo
PoC (import de SC do modelo, write-back de \texttt{VerificationResult}).

A distinção Requirement Table (SysML v1) vs.\ vistas SysML v2 no CATIA Magic 2026x
é relevante: a PoC não depende de CSV Sync como caminho principal.

\subsection{Telemetria e missões multiagente em laboratório}

O CONCEPTIO opera drones reais com publicação MQTT (JSON por entidade).
Trabalhos em arenas multiagente (ex.: WARA-PS) ilustram padronização de tópicos
e presença, mas não são meta de escala desta dissertação --- adoptam-se apenas
como referência de \emph{dialecto} futuro (\texttt{docs/MQTT\_PADRAO\_MENSAGEM.md}).

\section{Mapeamento da literatura}
\label{sec:slr}

\subsection{Protocolo de busca}

\begin{itemize}
  \item \textbf{Perguntas:} (P1) Como formalizar Success Criteria para medição
        live? (P2) Que gaps existem entre MBSE e telemetria real de UAS?
  \item \textbf{Fontes:} IEEE Xplore, Scopus, ACM; documentação NASA/ESA (MSFC,
        SIS-08); repositórios CONCEPTIO/CET-ADS.
  \item \textbf{Palavras-chave:} \emph{success criteria}, \emph{live verification},
        \emph{UAS telemetry}, \emph{SysML v2}, \emph{ground support equipment}, MQTT.
  \item \textbf{Inclusão:} trabalhos com critérios mensuráveis + evidência temporal;
        ligação explícita requisito--método--evidência.
  \item \textbf{Exclusão:} produtos comerciais sem artefacto replicável; simulação
        pura sem telemetria real quando o claim é V\&V de missão.
\end{itemize}

\subsection{Síntese}

\begin{table}[htbp]
  \centering
  \caption{Gap identificado vs.\ contribuição deste mestrado (síntese inicial).}
  \label{tab:slr-gap}
  \begin{tabular}{@{}lp{5.5cm}@{}}
    \toprule
    Tema na literatura & Lacuna / contribuição ReqValLive \\
    \midrule
    SC em MBSE & Critérios no modelo; gate antes de medir \\
    Telemetria UAS & Avaliação live multi-drone com latch temporal \\
    Laudos de V\&V & HTML auditável + snapshot SC congelado \\
    Write-back MBSE & \texttt{VerificationResult} no SysON (PoC) \\
    \bottomrule
  \end{tabular}
\end{table}

\section{Estado da arte do domínio do problema}
\label{sec:estado-arte}

\subsection{Naturezas de observação em GSE}

Contacto com especialista GSE (Jefferson / Sintra Lab) produziu taxonomia de
naturezas de observação~\cite{jefferson2026gse}: valor; tempo; sequência;
taxa; ocorrência/ausência; frequência; consistência entre estados.

A PoC implementa valor, amplitude em janela, taxa e dwell; as restantes alimentam
discussão e roadmap (\texttt{docs/JEFFERSON\_NATUREZAS\_OBSERVACAO.md}).

\subsection{Virtual Demonstrator e métricas cruzadas}

O ADS-VD / CET-ADS ancora a relevância industrial: medição cruzada entre
agentes reais e simulados. O ReqValLive posiciona-se na camada de \emph{validação
de requisitos declarados}, não na simulação física completa.
